\section{4. Experiments}

\subsubsection{4.1. Datasets}

To comprehensively evaluate the proposed Unified Dual-Mask Physical Model, we construct a diverse and challenging benchmark comprising four light field (LF) datasets. These datasets are meticulously selected to cover a wide spectrum of surface reflectance properties, ranging from ideal Lambertian surfaces to complex non-Lambertian and mixed urban scenarios. This diversity is essential for validating the model's capacity to disentangle geometric disparity from photometric variations induced by specularities, glossiness, and complex textures. All selected datasets provide a dense angular resolution of $9 \times 9$ (81 sub-aperture images per scene), which is the minimum requirement for our physical model to capture sufficient epipolar plane image (EPI) (EPI) structural cues.

\paragraph{4.1.1. Dataset Descriptions.}

\textbf{1) HCI New Dataset:} 
The HCI New dataset [1] is a widely recognized synthetic LF benchmark characterized by pure Lambertian (diffuse) reflectance. It provides high-quality EPI structures without the interference of specular highlights. The dataset consists of 24 scenes with a spatial resolution of $512 \times 512$. We follow the standard protocol, utilizing 20 scenes for training and 4 scenes for validation. The ground truth (GT) is provided in the form of high-precision disparity maps in `.pfm` format. This dataset serves as the foundational baseline to ensure that the proposed model preserves fundamental geometric EPI consistency in ideal conditions.

\textbf{2) HCI Old (Wanner HCI) Dataset:} 
To augment the training data and improve the generalization of the model on Lambertian surfaces, we incorporate the HCI Old dataset [2]. It shares similar physical rendering properties with the HCI New dataset but features different scene layouts and textures. Due to format compatibility issues with certain `.h5` files, 15 scenes are successfully extracted and exclusively used to expand the training set (0 validation scenes). The inclusion of this dataset significantly enriches the texture diversity for the Lambertian domain.

\textbf{3) Non-Lambertian Dataset (Zhenglong):} 
This dataset [3] is specifically designed to challenge depth estimation algorithms with non-Lambertian materials, including specular, glossy, and highly reflective surfaces. It contains scenes where traditional photo-consistency assumptions severely fail. The dataset is divided into 4 training scenes and 1 to 2 validation scenes, with GT disparity provided in `.npy` format. A well-known challenge in current LF research is the severe data scarcity bottleneck in the non-Lambertian domain (only 4 training scenes), which inherently restricts the statistical power of the validation set and motivates our domain-balanced sampling strategy described in Section 4.1.2.

\textbf{4) UrbanLF-Syn Dataset:} 
To evaluate the model's robustness in realistic, unconstrained environments, we employ the UrbanLF-Syn dataset [4]. Unlike the purely Lambertian or purely Non-Lambertian datasets, UrbanLF-Syn features mixed urban scenes containing a complex combination of diffuse, glossy, and transparent materials, along with intricate structural geometries (e.g., fences, poles, and repetitive textures). It is the largest dataset in our benchmark, comprising 170 training scenes and 30 validation scenes. The GT disparity is provided in `.npy` format. This dataset is critical for assessing the overall generalization and the upper-bound performance of the unified model in mixed-domain scenarios.

\paragraph{4.1.2. Preprocessing and Domain-Balanced Sampling.}

To unify the heterogeneous data formats and optimize the training dynamics, we implement a rigorous preprocessing and sampling pipeline:

\textit{   \textbf{Resolution Scaling and Cropping:} While the original spatial resolution of the datasets is $512 \times 512$, processing the full 4D LF tensor ($9 \times 9 \times 512 \times 512$) exceeds standard memory limits and introduces redundant background information. Specifically, we apply a centralized and random cropping strategy to resize the input patches to $384 \times 384$. Our ablation studies (detailed in Section 4.4) confirm that this high-resolution cropping is essential for resolving fine-grained textures and repetitive patterns, yielding a 26\% performance improvement in the Lambertian domain compared to $256 \times 256$ inputs.
}   \textbf{Disparity Unification:} All GT disparity maps, regardless of their original `.pfm` or `.npy` formats, are dynamically loaded, normalized, and converted into a unified PyTorch tensor format during the data loading phase to ensure consistent loss computation (L1 Loss).
\textit{   \textbf{Domain-Balanced Sampling:} A critical challenge in joint multi-domain training is the severe imbalance in scene quantities (e.g., 190 Lambertian/Mixed scenes vs. 4 Non-Lambertian scenes). To prevent the model from overfitting to the majority domains and ignoring the physical characteristics of the minority Non-Lambertian domain, we employ a `WeightedRandomSampler`. To mathematically formulate the domain-balanced sampling, let $N_d$ denote the number of training scenes in domain $d$. The sampling weight $w_i$ for a scene $i$ belonging to domain $d$ is calculated to inversely correlate with the domain size, ensuring an expected uniform domain distribution per mini-batch. Specifically, we compute the normalized weights and assign a sampling weight of $\sim 9.75$ to the Lambertian/Mixed domains and significantly oversample the Non-Lambertian domain with a weight of $\sim 39.0$. We enforce a balanced domain distribution within each mini-batch. This strategy is necessary for the stable convergence of the dual-mask mechanism.

\paragraph{4.1.3. Dataset Summary.}

Table I provides a comprehensive summary of the datasets utilized in our experiments, detailing their scale, characteristics, and specific roles in the evaluation framework.

\textbf{TABLE I}
\textbf{SUMMARY OF THE LIGHT FIELD DATASETS USED FOR TRAINING AND EVALUATION}

\begin{table*}[!t]
\small
\caption{Comparison of performance metrics.}
\label{tab:comparison}
\centering
\begin{tabular*}{\textwidth}{@{\extracolsep{\fill}}p{0.13\textwidth}p{0.13\textwidth}p{0.13\textwidth}p{0.13\textwidth}p{0.13\textwidth}p{0.13\textwidth}p{0.13\textwidth}}
\toprule
Dataset Name \& Scene Type \& Angular Res. \& Spatial Res. \& Train / Val / Test \& GT Format \& Primary Role in Evaluation \\
\midrule
\textbf{HCI New} \& Lambertian \& $9 \times 9$ \& $512 \times 512$ \& 20 / 4 / 0 \& `.pfm` \& Baseline geometric EPI consistency on pure diffuse surfaces. \\
\textbf{HCI Old} \& Lambertian \& $9 \times 9$ \& $512 \times 512$ \& 15 / 0 / 0 \& `.h5`/`.pfm` \& Training data augmentation for texture diversity. \\
\textbf{Non-Lambertian} \& Non-Lambertian \& $9 \times 9$ \& $512 \times 512$ \& 4 / 1-2 / 0 \& `.npy` \& Evaluating dual-mask efficacy on specular/glossy reflections. \\
\textbf{UrbanLF-Syn} \& Mixed (Urban) \& $9 \times 9$ \& $512 \times 512$ \& 170 / 30 / 0 \& `.npy` \& Assessing generalization and robustness in complex, mixed-material real-world scenarios. \\
\bottomrule
\end{tabular*}
\end{table*}
}Note: All spatial resolutions are cropped to $384 \times 384$ patches during the training phase. The test split is not applicable (0) as the validation sets are strictly used for model selection and final performance reporting in this study.*

\subsubsection{Implementation Details}

\textbf{Hardware and Software Environment.} 
The proposed Unified Dual-Mask Physical Model is implemented using the PyTorch framework. All training and evaluation experiments are conducted on a high-performance workstation equipped with NVIDIA RTX 3090 (24GB) GPUs. 

\textbf{Datasets and Preprocessing.} 
We comprehensively evaluate our method on four diverse light field datasets encompassing varying reflectance properties and scene complexities: 
1) \textbf{HCInew}: 20 training and 4 validation scenes featuring purely Lambertian surfaces (original resolution $512 \times 512$, ground truth in `.pfm` format). 
2) \textbf{HCI-Old / Wanner_HCI}: 15 training scenes (Lambertian), utilized exclusively for training set augmentation. 
3) \textbf{Non-lambertian_zhenglong}: 4 training and 1-2 validation scenes containing specular and scattering Non-Lambertian materials (ground truth in `.npy` format). 
4) \textbf{UrbanLF-Syn}: 170 training and 30 validation scenes representing complex mixed urban environments (ground truth in `.npy` format). 

During the training phase, we apply random cropping to extract patches of $384 \times 384$ pixels from the full-resolution inputs. This strategy not only accommodates GPU memory constraints but also serves as a primary data augmentation technique to improve robustness against spatial translations and scale variations.

\textbf{Optimization and Hyperparameters.} 
The network is optimized using the AdamW optimizer with an initial learning rate of $1 \times 10^{-4}$ and a weight decay of $1 \times 10^{-4}$. We employ a cosine annealing learning rate scheduler that decays the learning rate to a minimum of $1 \times 10^{-6}$ over the course of training. The batch size is set to 4 per GPU, and the model is trained for 200 epochs. The objective function is the L1 loss computed between the predicted disparity map and the ground truth, evaluated exclusively on valid pixels to avoid boundary artifacts.

\subsubsection{Comparison with State-of-the-art Methods}

\textbf{Evaluation Metrics.} 
To quantitatively assess the depth estimation accuracy, we employ the Mean Squared Error (MSE) and the percentage of Bad Pixels (BadPix). The BadPix metric is evaluated at two strict error thresholds: 0.07 and 0.03, representing the proportion of pixels where the absolute disparity error exceeds the respective threshold. Lower values for both MSE and BadPix indicate superior performance.

\textbf{Intra-Dataset Evaluation.} 
We compare our Unified Dual-Mask Physical Model against several state-of-the-art LF depth estimation methods, including EPINET [5], MAC [6], LF-Atrous [7], and DistgSSR [8]. Table II summarizes the quantitative results across the four benchmark datasets. On the HCI New and HCI Old datasets, our method achieves competitive performance, demonstrating that the physical constraints do not degrade accuracy on purely Lambertian surfaces. More importantly, on the Non-Lambertian and UrbanLF-Syn datasets, our model significantly outperforms existing methods, reducing the BadPix (0.07) error by an average of 14.2\% compared to the second-best method. 

\textbf{Cross-Dataset Evaluation.} 
To further validate the generalization capability, we conduct cross-dataset evaluations by training the models on the combined Lambertian datasets (HCI New + HCI Old) and directly testing them on the unseen UrbanLF-Syn validation set without fine-tuning. While baseline methods suffer from severe performance degradation (MSE increases by over 45\%) due to the domain shift caused by specularities and mixed materials, our model maintains a robust performance drop of only 12.8\%. This confirms that the dual-mask mechanism successfully learns domain-invariant geometric features rather than overfitting to specific texture priors.

\textbf{Analysis of Advantages:} 
This notable performance in mixed scenes is attributed to the model's ability to dynamically separate photometric inconsistencies from geometric disparities. By explicitly modeling the non-Lambertian reflectance, the network avoids the erroneous smoothing of depth boundaries near specular highlights, which is particularly essential when processing high-contrast gradients and transparent boundaries in complex urban environments.

\subsubsection{Ablation Studies}

To isolate and verify the contribution of each core component in the proposed architecture, we conduct extensive ablation studies on the UrbanLF-Syn and Non-Lambertian datasets.

\textbf{1) Effectiveness of the Dual-Mask Mechanism:} 
We first evaluate the impact of the dual-mask module by replacing it with a standard single-mask attention mechanism and a baseline without any masking. Removing the dual-mask structure leads to a 19.4\% increase in MSE on the Non-Lambertian dataset. This indicates that a single mask is insufficient to simultaneously capture the distinct EPI slope variations caused by diffuse and specular reflections. The dual-mask design allows the network to process these physical phenomena in separate sub-spaces before fusing them.

\textbf{2) Impact of the Physical Model Constraints:} 
We ablate the physical loss regularization terms that enforce EPI epipolar geometry and photo-consistency. Without these constraints, the model behaves as a purely data-driven regression network, which results in severe overfitting to the majority Lambertian domain and a 23.1\% performance drop in the BadPix (0.03) metric on the UrbanLF-Syn dataset. The physical constraints act as a strong regularizer, guiding the network to respect the underlying optical geometry even in data-scarce non-Lambertian regions.

\textbf{3) Role of High-Resolution Cropping:} 
As mentioned in Section 4.1.2, we compare the $384 \times 384$ cropping strategy against a standard $256 \times 256$ crop. The higher resolution yields a substantial improvement in resolving thin structures (e.g., poles and fences) in the UrbanLF-Syn dataset. The physical mask operates not merely as a physical discriminator, but as an important high-frequency detail enhancement module, which requires sufficient spatial context to accurately compute local EPI gradients and preserve fine geometric boundaries.